\documentclass[12pt,parskip=half-, twoside=semi]{scrartcl}
\usepackage{mathtools}
\usepackage[T1]{fontenc}
\usepackage[utf8]{inputenc}
\usepackage{multicol}
\usepackage{cmbright}  %police sans serif
\usepackage{enumitem}
\SetEnumitemKey{cols}
{
  before = \begin{multicols}{#1},
    after = \end{multicols}
    }
\setlist[enumerate]{label=\textbf{\alph*)}, itemsep=1em}
\usepackage[use-files]{xsim}


\DeclareExerciseEnvironmentTemplate{mytemplate}
{%
    \medbreak
    {\bfseries\XSIMmixedcase{\GetExerciseName}\nobreakspace\GetExerciseProperty{counter}.}%
    \GetExercisePropertyT{points}{\marginpar{\makebox[4ex]{\dotfill}/\bfseries\PropertyValue}}
    \GetExercisePropertyT{subtitle}{\normalfont\itshape\PropertyValue\par}
    }%
{\par}

\xsimsetup{
    path=xsim,
    exercise/template=mytemplate
}


\usepackage[french]{babel}
\usepackage{fontawesome5}
\usepackage{tabularray}
\usepackage{scrlayer-scrpage}
%  \setlength{\headheight}{50pt}
%  \lohead{Fiches d'exercices de drill}
%  \rohead{Chapitre 4 - Fractions}
\pagestyle{empty}


\begin{document}

%Niveau 1: une opération +/-/x/:
%Niveau 2: opération dans les 2 membres
%Niveau 3: avec fractions
%Niveau 4: dépassement

\begin{center}
  \bfseries
  \begin{Huge}
    Chapitre 9 - Equations
  \end{Huge}
  \smallbreak
  \begin{large}
    Niveau 1 \faIcon{star} - Fiche 1
  \end{large}
  \smallbreak
  \rule[10pt]{.5\linewidth}{1 pt}
\end{center}
\setcounter{exercise}{0}

\begin{exercise}
  Résous les équations suivantes.
  \begin{enumerate}[itemsep=0.4cm]
    \item \(x + 3 = 6\)
    \item \(2=x-4\)
    \item \(-x + 5 = -6\)
    \item \(x - \dfrac{1}{2} = 10\)
    \item \(-8=x + 3\)
    \item \(10=9-(-2x)\)
    \item \(\dfrac{3}{2}+3x=2\)
    \item \(-16=-5x-1\)
    \item \(5x-3=-23\)
    \item \(7=9-2x\)
    \item \(\dfrac{10}{2} x  = -10\)
    \item \(11- x = -6\)
\end{enumerate}
\end{exercise}

\begin{center}
    \bfseries
    \begin{Huge}
      Chapitre 9 - Equations
    \end{Huge}
    \smallbreak
    \begin{large}
      Niveau 1 \faIcon{star} - Fiche 2
    \end{large}
    \smallbreak
    \rule[10pt]{.5\linewidth}{1 pt}
  \end{center}
  \setcounter{exercise}{0}

\begin{exercise}
    Résous les équations suivantes.
    \begin{enumerate}[itemsep=0.4cm]
    \item \(\dfrac{9}{4} x+9=8\)
    \item \(-4=x - 1\)
    \item \(\dfrac{3}{4} x=7\)
    \item \(9x=27\)
    \item \(-3x=9\)
    \item \(9=6+3x\)
    \item \(3+4x=19\)
    \item \(5=2x-12\)
    \item \(-\dfrac{1}{4} x = 8\)
    \item \(-x+32=27\)
    \item \(2x+9=21\)
    \item \(12=-5x+ 5\)
  \end{enumerate}
\end{exercise}

\begin{center}
    \bfseries
    \begin{Huge}
      Chapitre 9 - Equations
    \end{Huge}
    \smallbreak
    \begin{large}
      Niveau 1 \faIcon{star} - Fiche 3
    \end{large}
    \smallbreak
    \rule[10pt]{.5\linewidth}{1 pt}
  \end{center}
  \setcounter{exercise}{0}

\begin{exercise}
    Resolve the following equations.
    \begin{enumerate}[itemsep=0.4cm]
        \item \( \dfrac{5}{2}x - 3 = 7 \)
        \item \( 2x + 1 = 5 \)
        \item \( \dfrac{3}{5}x = 10 \)
        \item \( 4x = 16 \)
        \item \( -2x = 6 \)
        \item \( 7 = 3x + 4 \)
        \item \( 2x - 8 = 14 \)
        \item \( 6 = \dfrac{1}{3}x + 2 \)
        \item \( -\dfrac{2}{7}x = 4 \)
        \item \( x + 5 = 10 \)
        \item \( 3x - 2 = 13 \)
        \item \( 9 = -2x + 7 \)
    \end{enumerate}
    \end{exercise}

    \begin{center}
        \bfseries
        \begin{Huge}
          Chapitre 9 - Equations
        \end{Huge}
        \smallbreak
        \begin{large}
          Niveau 1 \faIcon{star} - Fiche 4
        \end{large}
        \smallbreak
        \rule[10pt]{.5\linewidth}{1 pt}
      \end{center}
      \setcounter{exercise}{0}
    \begin{exercise}
        Resolve the following equations.
        \begin{enumerate}[itemsep=0.4cm]
            \item \( 3x - 4 = 10 \)
            \item \( \dfrac{2}{5}x + 3 = 7 \)
            \item \( \dfrac{1}{2}x = 6 \)
            \item \( 5x = 25 \)
            \item \( -2x = -8 \)
            \item \( 8 = 2x + 4 \)
            \item \( x + 6 = 18 \)
            \item \( 9 = \dfrac{1}{4}x + 5 \)
            \item \( -\dfrac{3}{2}x = 12 \)
            \item \( 2x - 7 = 5 \)
            \item \( 4x + 3 = 15 \)
            \item \( 6 = -3x + 9 \)
        \end{enumerate}
        \end{exercise}

\begin{center}
    \bfseries
    \begin{Huge}
      Chapitre 9 - Equations
    \end{Huge}
    \smallbreak
    \begin{large}
      Niveau 2 \faIcon{star}\faIcon{star} - Fiche 1
    \end{large}
    \smallbreak
    \rule[10pt]{.5\linewidth}{1 pt}
  \end{center}
  \setcounter{exercise}{0}

\begin{exercise}
  Résous les équations suivantes.
    \begin{enumerate}[itemsep=0.4cm]
    \item \(\dfrac{1}{2}+3x=\dfrac{4}{3}\)
    \item \(55x-15=75x+18\)
    \item \(2x+\dfrac{7}{3}=\dfrac{1}{4}\)
    \item \(4x+3=2x-7\)
    \item \(x-1=3x+9\)
    \item \(2x+8\cdot 5-5=7\cdot 3x\)
    \item \(-14+5x\cdot 2=x+5\)
    \item \(\dfrac{2x}{3}-\dfrac{5x}{4}+5=-6x+\dfrac{4}{3}\)
    \item \(12x+4=88x-16\)
    \item \(-8-8x+2=-6x+x\)
    \item \(19+3x-5x+1=11x\)
    \item \(\dfrac{-5}{2}+\dfrac{2x}{3}+x=\dfrac{4}{3}\)
    \item $20 = 8x + 2x$
    \item $6x = \dfrac{6}{3}$
    \item $- 3x + 5 = - 1$
    \item $(3 + 2x)^{2} = 5x + 4x^2$
    \end{enumerate}
\end{exercise}

\begin{center}
    \bfseries
    \begin{Huge}
      Chapitre 9 - Equations
    \end{Huge}
    \smallbreak
    \begin{large}
      Niveau 2 \faIcon{star}\faIcon{star} - Fiche 2
    \end{large}
    \smallbreak
    \rule[10pt]{.5\linewidth}{1 pt}
  \end{center}
  \setcounter{exercise}{0}

\begin{exercise}
    Résous les équations suivantes.
    \begin{enumerate}[itemsep=0.4cm]
    \item $3(x + 3) = 3x + 9$
    \item $\dfrac{12}{x} + 3 = 6$
    \item $7(5 - 2x)+5 = 45$
    \item $- 7x = \dfrac{5}{6}$
    \item $(x - 9)(1 + x) - 7 = x^2$
    \item $- (5 - x) + 2x = 5x$
    \item \(-7x+\dfrac{21}{3}=-43+2x\)
    \item \(\dfrac{6x}{7}+0,2=-x-\dfrac{3}{21}  \)
    \item \( 2x+\dfrac{7}{3}=\dfrac{1}{4}\)
    \item \(9-2x=7x+1\)
    \item \(-125x-25x=80-10x \)
    \item \(\dfrac{-8x}{15}+\dfrac{2}{5}=0,25x+7\)
    \item \(-36+56-2x=18x\)
    \item \(-18+45x+2x=50x\)
    \item \(-14+5x\cdot 2=x+5\)
  \end{enumerate}
\end{exercise}

\begin{center}
    \bfseries
    \begin{Huge}
      Chapitre 9 - Equations
    \end{Huge}
    \smallbreak
    \begin{large}
      Niveau 2 \faIcon{star}\faIcon{star} - Fiche 3
    \end{large}
    \smallbreak
    \rule[10pt]{.5\linewidth}{1 pt}
  \end{center}
  \setcounter{exercise}{0}

\begin{exercise}
    Résous les équations suivantes.
    \begin{enumerate}[itemsep=0.4cm]
    \item \( 3(x + 3) = 3x + 9 \)
    \item \( \dfrac{12}{x} + 3 = 6 \)
    \item \( 7(5 - 2x)+5 = 45 \)
    \item \( - 7x = \dfrac{5}{6} \)
    \item \( (x - 9)(1 + x) - 7 = x^2 \)
    \item \( - (5 - x) + 2x = 5x \)
    \item \( -7x+\dfrac{21}{3}=-43+2x \)
    \item \( \dfrac{6x}{7}+0,2=-x-\dfrac{3}{21}  \)
    \item \( 2x+\dfrac{7}{3}=\dfrac{1}{4} \)
    \item \( 9-2x=7x+1 \)
    \item \( -125x-25x=80-10x \)
    \item \( \dfrac{-8x}{15}+\dfrac{2}{5}=0,25x+7 \)
    \item \( -36+56-2x=18x \)
    \item \( -18+45x+2x=50x \)
    \item \( -14+5x\cdot 2=x+5 \)
  \end{enumerate}
\end{exercise}

\begin{center}
    \bfseries
    \begin{Huge}
      Chapitre 9 - Equations
    \end{Huge}
    \smallbreak
    \begin{large}
      Niveau 2 \faIcon{star}\faIcon{star} - Fiche 4
    \end{large}
    \smallbreak
    \rule[10pt]{.5\linewidth}{1 pt}
  \end{center}
  \setcounter{exercise}{0}

\begin{exercise}
    Résous les équations suivantes.
      \begin{enumerate}[itemsep=0.4cm]
      \item \( \dfrac{1}{2}+3x=\dfrac{4}{3} \)
      \item \( 55x-15=75x+18 \)
      \item \( 2x+\dfrac{7}{3}=\dfrac{1}{4} \)
      \item \( 4x+3=2x-7 \)
      \item \( x-1=3x+9 \)
      \item \( 2x+8\cdot 5-5=7\cdot 3x \)
      \item \( -14+5x\cdot 2=x+5 \)
      \item \( \dfrac{2x}{3}-\dfrac{5x}{4}+5=-6x+\dfrac{4}{3} \)
      \item \( 12x+4=88x-16 \)
      \item \( -8-8x+2=-6x+x \)
      \item \( 19+3x-5x+1=11x \)
      \item \( \dfrac{-5}{2}+\dfrac{2x}{3}+x=\dfrac{4}{3} \)
      \item \( 20 = 8x + 2x \)
      \item \( 6x = \dfrac{6}{3} \)
      \item \( - 3x + 5 = - 1 \)
      \item \( (3 + 2x)^{2} = 5x + 4x^2 \)
      \end{enumerate}
  \end{exercise}


  \begin{center}
    \bfseries
    \begin{Huge}
      Chapitre 9 - Equations
    \end{Huge}
    \smallbreak
    \begin{large}
      Niveau 3 \faIcon{star}\faIcon{star}\faIcon{star} - Fiche 1
    \end{large}
    \smallbreak
    \rule[10pt]{.5\linewidth}{1 pt}
  \end{center}
  \setcounter{exercise}{0}

  %Equations algébriques
  \begin{exercise}
    Résous les équations algébriques suivantes.
    \begin{enumerate}[itemsep=0.4cm]
    \item $\dfrac{4 x+3}{2}=\dfrac{5-2 x}{6}$
    \item $\dfrac{2 x+5}{6}-\dfrac{x-8}{4}=\dfrac{3-2 x}{2}-\dfrac{6-4 x}{3}$
    \item $\dfrac{-3 x+9}{4}=\dfrac{7 x+2}{3}$
    \item $-5+\dfrac{7-9 x}{2}=\dfrac{1}{5}-\dfrac{4 x+3}{10}$
    \item $\dfrac{10 x-5}{2}=\dfrac{-6+x}{10}$
    \item $\dfrac{-x}{4}+\dfrac{-2 x+4}{3}+2=\dfrac{5 x-1}{6}+3 x$
\end{enumerate}
\end{exercise}

\begin{exercise}
    Résous les équations suivantes.
      \begin{enumerate}[itemsep=0.4cm]
    \item $\dfrac{-3-8 x}{7}=\dfrac{9+2 x}{2}$
    \item $\dfrac{1}{2}-\dfrac{6+3 x}{5}=\dfrac{9}{2}+x$
    \item $\dfrac{7-9 x}{24}=\dfrac{8+x}{8}$
    \item $-\dfrac{2-7 x}{6}+\dfrac{3-4 x}{5}+x=\dfrac{9 x+3}{10}+\dfrac{1}{5}$
    \item $\dfrac{11 x+1}{33}=\dfrac{-x-2}{22}$
    \item $\dfrac{15+x}{3}+\dfrac{2 x}{3}=\dfrac{9}{5}$
\end{enumerate}
\end{exercise}

\begin{exercise}
    Résous les équations suivantes.
      \begin{enumerate}[itemsep=0.4cm]
    \item $\dfrac{6+2 x}{8}=\dfrac{10-x}{12}$
    \item $\dfrac{-6+x}{7}-\dfrac{10}{21}=\dfrac{5 x+2}{3}+\dfrac{3}{7}$
    \item $\dfrac{-7 x-6}{4}=\dfrac{2 x+1}{5}$
    \item $\dfrac{-5}{2}+\dfrac{x-9}{3}-\dfrac{8 x+8}{12}=5$
    \item $\dfrac{-x-1}{17}=\dfrac{1+x}{51}$
    \item $\dfrac{11 x+10}{10}=\dfrac{1}{5}+5 x+\dfrac{4-8 x}{20}$
\end{enumerate}
\end{exercise}

\begin{exercise}
    Résous les équations suivantes.
      \begin{enumerate}[itemsep=0.4cm]
    \item $\dfrac{3 x+2}{16}=\dfrac{-x}{48}$
    \item $8+\dfrac{1+2 x}{4}=\dfrac{6 x+3}{5}-\dfrac{1}{2}+\dfrac{x-9}{4}$
    \item \( \dfrac{-7x-6}{4} = \dfrac{2x+1}{5} \)
    \item \( \dfrac{-5}{2} + \dfrac{x-9}{3} - \dfrac{8x+8}{12} = 5 \)
    \item \( \dfrac{-x-1}{17} = \dfrac{1+x}{51} \)
    \item \( \dfrac{11x+10}{10} = \dfrac{1}{5} + 5x + \dfrac{4-8x}{20} \)
\end{enumerate}
\end{exercise}
  
    \begin{center}
    \bfseries
    \begin{Huge}
      Chapitre 9 - Equations
    \end{Huge}
    \smallbreak
    \begin{large}
      Niveau 4 \faIcon{crown} - Fiche 1
    \end{large}
    \smallbreak
    \rule[10pt]{.5\linewidth}{1 pt}
  \end{center}
  \setcounter{exercise}{0}

\begin{exercise}
  Résous les équations suivantes.
    \begin{enumerate}[itemsep=0.4cm]
        \item $5 x+3 \cdot(-4)+12=(4 x+3) \cdot 2$
        \item $5 x+8-9=2 x+7$
        \item $7+\dfrac{2 x}{3}=8-\dfrac{5}{2} x$
        \item $\dfrac{1}{2}-(4 x+6)=-3 x+2$
        \item $-(4 x+2)^2+5=4 x \cdot(-4-4 x)+1$
        \item $\dfrac{3 x}{\dfrac{9}{5}}+x \cdot(-3)=x$
    \end{enumerate}
\end{exercise}

\begin{exercise}
    Résous les équations suivantes.
      \begin{enumerate}[itemsep=0.4cm]
        \item $\dfrac{7 x}{\dfrac{-2}{3}}+9 x=\dfrac{5}{4}+\dfrac{3 x}{5}$
        \item $-(4 x \cdot 3+5-2 x)=-2 \cdot(3 x+1)$
        \item $2 x+9:(-3)=4 \cdot(-2 x-8)+x$
        \item $\dfrac{-5 x}{3}+2 x=\dfrac{-8}{7}-2$
        \item $(5 x-1) \cdot(1+5 x)=x+5 x \cdot 5 x$
        \item $\left(\dfrac{2}{5}\right)^{-1}+x=-\dfrac{-3}{10}+2 x$
    \end{enumerate}
\end{exercise}

\begin{exercise}
    Résous les équations suivantes.
      \begin{enumerate}[itemsep=0.4cm]
        \item $\dfrac{1}{6} \cdot \dfrac{-3 x}{4}+\dfrac{7}{16}=3$
        \item $(x-1) \cdot(x+1)+8=3+x \cdot x-x$
        \item $-[3 x+2:(-2)+10]=5 x+3$
        \item $\dfrac{2}{4 x-3}=\dfrac{9}{2 x}$
        \item $\dfrac{\dfrac{2 x}{3}}{2}=5 x+1$
        \item $(x+8) \cdot 2-(5 x-1)=7 x+12$
    \end{enumerate}
\end{exercise}

\begin{exercise}
    Résous les équations suivantes.
      \begin{enumerate}[itemsep=0.4cm]
        \item $\dfrac{15}{-25}+2 \cdot(x+4 x-11 x)=-(3 x+5-9 x)$
        \item $\dfrac{\dfrac{15 x}{3}}{\dfrac{5}{3}}+4 x-2=12:(-6) \cdot(-3)+4 x$
        \item $18 x+4 \cdot 10 x:(-8)+1=7 x-6+\dfrac{125}{25}$
        \item $\dfrac{144}{-12}+5 x=\dfrac{-4}{2}-x$
        \item $6 \cdot(15 x: 3)+7=6 \cdot 5 x-7$
        \item $\dfrac{44 x}{11}+8=-(5 x+6-10 x)+1$
        
        \item $2 x \cdot(-5+9)=\dfrac{51}{17}+3 x$
        \item $-8-\dfrac{16 x}{4}+4 x=2 \cdot(-4 x)+5$
    \end{enumerate}
\end{exercise}

\end{document}